\section{Knowledge Gap}
There are two main adverse effects of radar receivers that are sized for small to medium UAVs. The first is low received power, which is proportional to the effective receive aperture area. Antenna gain \(G\) scales with effective aperture \(A_e\) as \(G \propto A_e\). All else equal, reduced gain lowers received power and therefore reduces SNR, with \(\mathrm{SNR} \propto G^2\), decreasing the sensitivity of the system to small or far away targets. The second adverse effect is poor angular resolution. Beamwidth \(\theta_{\mathrm{BW}}\) scales as \(\theta_{\mathrm{BW}} \propto 1/D\), where \(D\) is the effective aperture diameter. With a small radar system and thus a small effective diameter, beamwidth is large and resolution is poor. Therefore, it becomes more difficult to resolve the velocity and positions of single targets, differentiate between multiple targets, identify their size, and have confidence in detections. 

\todo{Talk broadly about the state of the current literature and how they engage with this new issue. Summarise the major gaps in the literature in terms of radar systems under swap constraints, and if the literature has missed any techniques that we're going to employ. Tie that into our aims and outcomes, explaining how our experiments will fill the gaps.}